\chapter{Against Homura}
There is a recurring image throughout the Madoka Magica movies, one we have briefly mentioned before: a rather sweet tableau of two white chairs on a grassy hill, Madoka and Homura sitting side-by-side in them. In the opening credits of the first two movies, they cuddle, sweet and adorable, and innocent. In the third movie, the image turns rapidly rather less sweet.

As she goes through the process of becoming a witch at the climax of the second arc of Rebellion, Homura returns to the chair scene. But this time, Madoka stands and casts herself sideways off the chair, splattering into a pink stain on the grass while Homura reaches for her helplessly. Homura crouches beside her, eyes wide in shock and horror, while a crowd of tall, attenuated Homuras surround her, gazing down. And then the vast fist of a raging Homura smashes the crouching Homura, railing and weeping beside the remains of Madoka.

Madoka is gone, her coherent identity replaced by a diffuse abstraction. Homura failed. Now Homura stands in judgment over Homura, and finds her wanting. Her rage and grief at last unleashed, she smashes her own identity to become an abstract and esoteric being herself: a witch. 

Just like Sayaka, and presumably every other witch before her, Homura's witch form is an endless cycle of self-flagellation, a psychodrama in which she acts out the events that brought her to despair and punishes herself for her failures. She tries to shoot herself, and the self she shoots becomes the Madoka she had to mercy-kill. She cannot die, does not deserve to die the way that Madoka did, because she has failed to save Madoka.

Not only failed to save her; Homura is the reason Madoka is gone. Her looping through time empowered Madoka to become the Law of Cycles, which erased Madoka from reality. Her discussion of Madoka with Kyubey gave the Incubators the information they needed to construct the trap now closing on Madoka—and they used Homura to create that trap. Homura is Madoka's greatest liability.

Homura's witch form is among the most literal. She has the peaked black hat, the prominent nose and chin—other than being a skeleton hundreds of feet tall, she looks rather like the standard Halloween costume of a witch. Homura knew about witches and where they come from, and yet she still failed to avoid that trap, even embraced it deliberately in a bid to foil Kyubey. Unlike Sayaka, who believed herself a knight and so still looked like one as a witch, Homura knows what she is choosing to become. Likewise, she is deliberately sacrificing herself, as she tells Kyubey: she trusts Mami and Kyoko to kill her. Thus her familiars lead her to the guillotine, the mechanism of her sacrifice and instrument of judgment for her crime.

At the same time, she is surrounded by imagery related to the nutcracker. One type of her familiars is giant teeth with nutcracker jaws. Another resembles toy soldiers, but with their high fur hats resemble the traditional Christmas nutcracker as well (and in the original story "The Nutcracker and the Mouse King," the nutcracker is immediately recognized as a soldier by the children \cite{ref127}). An image of a grinning mouth clenching a walnut in its teeth appears when she first starts to realize that she is the witch in whose labyrinth the magical girls are trapped. And she loses half her head, leaving only the lower jaw—a mirror of the titular nutcracker of E.T.A. Hoffman's story and Tchaikovsky's famous ballet based on it, who lost his lower jaw. The doll-like appearance of many of her familiars (in particular, the ragdolls, named "Clara" in supplementary materials \cite{ref128}, the name of the main character in the ballet), the image of Homura and her familiars breaking out of a glass-fronted cabinet, and prominence of clockwork also recall the original story of "The Nutcracker and the Mouse King," in which the Nutcracker led an army of toys from a cabinet against the Mouse King's armies summoned by the striking of a clock. 

At a basic level, the image of a nutcracker without a jaw is an image of uselessness, an object without purpose. There is a deeper resonance here, however, if one recalls the tale-within-a-tale of the origin of the nutcracker in Hoffman's story. The nutcracker was once the chosen one, described in what amounts to prophecy as the only one who could rescue a princess cursed by the Mouse Queen. He had to perform a complex ritual to save her, but just as he completed it, he tripped over the Mouse Queen, and so the curse fell on him instead. This is Homura, relaxing because she believed she had helped Madoka escape her fate, only to discover that she'd failed in the end because of the intervention of Kyubey. It is, in other words, yet another way to blame and punish herself.

Yet the magical girls refuse to cooperate. They refuse to join Homura in judging herself. They refuse to hate her and refuse to kill her. Instead, they work to free her, break the labyrinth and the Incubators' trap so that Madoka can take her off to magical girl heaven. Despite her raving and her pleading, they insist on forgiving her. They reject Homura's judgment, and demand that she reject it as well. They want her to forgive herself and free herself. 

But Homura has been fighting Homura from the start of the movie. Throughout the first arc of the film, Homura seeks the mysterious and invisible tyrant who rules the seemingly happy world in which the magical girls find themselves, with the intent of destroying it. It is the discovery that she is that tyrant which leads her to call down a curse on herself and transform fully into a witch; all of this is part of her rebellion against herself. 

That rebellion has not ended by the end of the film. Homura describes herself as evil and embraces the role of the scantily clad, black-winged devil-woman. But what difference is there between saying "I am evil," and "I deserve to be punished?" This is simply another expression of her guilt, a new way of tormenting herself. 

She has elevated herself to a cosmic being, a demiurgic entity who appears to have near-unlimited powers over material reality and the people in it: she can rewrite Sayaka's memories, bring back the dead, construct an entire new history for Madoka's family in order to reverse the first episode. And yet she chooses to make a world where she is alone, isolated from the friendships she was starting to build with the other magical girls. She chooses to let Sayaka tell her off before the memory erasure.

The only real emotion Homura shows in the new reality she created is panic, when Madoka threatens to reconnect with the Law of Cycles. When, in other words, Madoka nearly brings about the return of a cosmic entity of hope and forgiveness, capable of ending Homura's suffering. Above all, Homura cannot allow that; she must suffer for failing Madoka, making things worse for Madoka. She must preserve Madoka eternally in a state of innocence and safety, cut off from her potential, because protecting Madoka is Homura's only concept of "good"—and so her failure to do so is her only concept of "evil."

It could have ended. If the other magical girls had simply killed her, she would be beyond further punishment, and her suffering would have ended. But they, in their cruel mercy, forced her to go on, forced her to find another way to keep protecting Madoka and punishing herself. She hates them for that, for failing to hate her as she hates herself. In her new world, she expresses her hatred by passive-aggressively mocking its targets. She breaks a teacup behind Mami, recalling her death fighting Charlotte. She taunts Sayaka as her memories decay, mimicking Sayaka's loss of self when she became a witch. She tricks Kyoko into wasting food.

And, in the stinger, she throws herself off a cliff next to a white chair, mirroring Madoka tipping off it earlier. Her hatred for herself has not changed. All that has changed is that now she has the power to make the magical girls hate her, to position herself as their enemy in the hopes that they will finish the job.

Ever since the movie aired, there has been debate over Homura's new status. Is she hero or villain? Here, then, is the answer to that question: Yes. Homura is both the villain of Rebellion and the hero battling that villain.

And here, also, is the answer to that question: No. Homura is the villain's victim, whom the hero must rescue.

Her witch's barrier expanded to encompass the universe. She is the entire story now.
